% =============================================================================
% 1.2 The Research Problem
% =============================================================================

\section{The Research Problem}
\label{sec:research-problem}

\subsection{The Surface Problem}
\label{subsec:surface-problem}

A VM-hosted workflow orchestration system becomes unreliable under concurrent workload because all jobs share the same compute resources. This is an observable and known failure mode.

\subsection{The Deeper Problem --- The Research Gap}
\label{subsec:research-gap}

The surface problem is not the research gap. The gap exists at the intersection of two bodies of literature that have not yet been connected.

Existing research evaluates Kubernetes performance and autoscaling behaviour at the infrastructure level---measuring container overhead, pod scheduling latency, and cluster autoscaler response under generic workloads \citep{casalicchio2019, nguyen2020, tran2022}. Separately, research on resource contention analyses performance degradation in shared execution environments \citep{nanda1991, arora2023}. What is missing is empirical measurement of how infrastructure-level Kubernetes mechanisms---pod isolation, resource limits, and horizontal autoscaling---actually affect application-level workflow execution behaviour under controlled concurrent workload conditions.

\citet{casalicchio2019} demonstrates that under concurrent workloads with CPU utilisation above 75\%, response times under the default Kubernetes Horizontal Pod Autoscaler are 2 to 3 orders of magnitude higher than under a contention-aware algorithm, highlighting that resource contention remains severe even in containerised environments. \citet{choi2021} report that the median container startup latency in production orchestration systems is 25 seconds, representing a substantial fixed overhead for every pipeline job. But the specific workload threshold at which migrating from VM-based to Kubernetes-based workflow execution changes the balance between contention cost and scheduling overhead has not been empirically identified.

This thesis bridges that gap by measuring system behaviour at the pipeline execution level rather than at the infrastructure level alone.

\subsection{The Literature Gap Statement}
\label{subsec:gap-statement}

Existing studies focus on Kubernetes performance and autoscaling at the infrastructure level, but there is no empirical evidence on how these mechanisms affect application-level workflow orchestration systems under controlled concurrent workload conditions. No study identifies the specific concurrency threshold at which a Kubernetes migration becomes net beneficial for pipeline orchestration. This thesis addresses that specific gap.
